\chapter{Critical completion and constructible coefficients}
\label{vc:chap:critical}

The incidence equations of Chapter~\ref{ei:chap:incidence}
define a coherent algebra. A vanishing-cycle coefficient
requires a potential. For the triangle, three equations
constrain two coordinates, so the given derived zero locus
cannot itself carry a nondegenerate $(-1)$-shifted form.
One multiplier for each equation produces a critical space.
Its additional directions determine the coefficient to be
compared with compact Hall geometry.

We first construct this completion, including its coordinate
changes, orientation and constraint correspondences.
The function $\lambda z^m$ then separates three kinds of
data: coherent multiplicity, transverse vanishing monodromy
and compact integration. For the triangle, we determine
the entire rational central perverse sheaf and its two
monodromy summands. One nonzero summand has zero compact
integral. The displacement family also retains a nontrivial
specialization map when nearby monodromy is the identity.

All schemes are over $\mathbb C$. Differentials have degree $1$,
and $C[r]^i=C^{i+r}$. Fix a smooth separated scheme $B$
of finite type and pure dimension $m$,
a rank-$r$ vector bundle $V\to B$, and a section $s$ of $V$.
Its derived zero locus $Z=Z^{\mathrm{der}}(s)$ has the Koszul
algebra $\operatorname{Sym}_{\mathcal O_B}(V^\vee[1])$,
with differential given by contraction with $s$.
Write $Z_{\mathrm{cl}}=t_0Z$ for its classical truncation.

The cotangent complex $\mathbb L_Z$ represents derived
first-order differentials. Its dual is
\[
 \mathbb T_Z=R\mathcal Hom(\mathbb L_Z,\mathcal O_Z).
\]
A nondegenerate $(-1)$-shifted two-form induces a
quasi-isomorphism $\mathbb T_Z\simeq\mathbb L_Z[-1]$.

\section{The tangent complex and the critical space}

\begin{proposition}\label{vc:cicp:prop:rank}
The cotangent complex of $Z$, restricted to $Z_{\mathrm{cl}}$,
has the two-term presentation
\begin{equation}\label{vc:cicp:eq:cotangent}
 \mathbb L_Z|_{Z_{\mathrm{cl}}}
   =\bigl[V^\vee|_{Z_{\mathrm{cl}}}
       \xrightarrow{\,ds\,}\Omega_B|_{Z_{\mathrm{cl}}}\bigr]
 \quad\text{in degrees }-1,0.
\end{equation}
If $Z_{\mathrm{cl}}$ is nonempty and $Z$ has a nondegenerate
$(-1)$-shifted two-form, then $m=r$.
In particular, the fixed-parameter derived incidence scheme
of a connected graph with a cycle, using the pair sections
$\vartheta(u_e)\vartheta(u_e+h)$, has no such form.
\end{proposition}

\begin{proof}
The free Koszul presentation gives \eqref{vc:cicp:eq:cotangent}.
At any classical point its Euler rank is $m-r$.
Nondegeneracy would identify
$\mathbb T_Z$ with $\mathbb L_Z[-1]$. Duality preserves
Euler rank and this shift reverses it, so $m-r=-(m-r)$.

For a connected graph on $n$ vertices and $e$ edges, quotienting
simultaneous spectral translation leaves $m=n-1$ coordinates
and $r=e$ equations. The incidence locus contains the
configuration with every difference zero. Thus $m-r=-\beta$, where
$\beta=e-n+1$ is the number of independent cycles.
\end{proof}

Nonemptiness is necessary: the section $1$ of the trivial
line over $\mathbb A^2$ has empty derived zero locus, since
its Koszul algebra has $d\epsilon=1$. Nondegeneracy on the
empty scheme imposes no rank equality.

This calculation concerns the specified derived enhancement.
It permits another critical enhancement of the same classical
support. It also permits an absolute parameter family with a
different Euler rank. For a displacement parameter $h$, the
assertion applies fiberwise, or to the relative cotangent
complex over the $h$-curve. Including $h$ among the coordinates
adds one to the rank calculation.

The shifted cotangent space is
$T^*[-1]Z=\operatorname{Spec}_Z
\operatorname{Sym}_{\mathcal O_Z}(\mathbb T_Z[1])$.

Let $M=\operatorname{Tot}_B(V^\vee)$, with projection $\pi$.
The evaluation pairing defines the regular function
\begin{equation}\label{vc:cicp:eq:potential}
 F(x,\lambda)=\langle\lambda,s(x)\rangle.
\end{equation}
Its derived critical locus, denoted $\widehat Z$, is the derived
zero locus of $dF$ on $M$.

\begin{theorem}\label{vc:cicp:thm:completion}
There is an equivalence of $(-1)$-shifted symplectic derived
schemes
\begin{equation}\label{vc:cicp:eq:completion}
 \widehat Z=\operatorname{Crit}^{\mathrm{der}}(F)
       \simeq T^*[-1]Z.
\end{equation}
In local coordinates $x_1,\ldots,x_m$ and a frame
$s_1,\ldots,s_r$, its algebra is
\begin{equation}\label{vc:cicp:eq:critical-algebra}
 \widehat A=\mathcal O_B[\lambda_1,\ldots,\lambda_r]
       \otimes\Lambda(\xi_1,\ldots,\xi_m,
                       \epsilon_1,\ldots,\epsilon_r),
\end{equation}
where $|\lambda_a|=0$, $|\xi_i|=|\epsilon_a|=-1$, and
\begin{equation}\label{vc:cicp:eq:differential}
 d\epsilon_a=s_a,\qquad
 d\xi_i=\sum_a\lambda_a\partial_i s_a,\qquad
 dx_i=d\lambda_a=0.
\end{equation}
The inclusion of the original Koszul algebra represents
$T^*[-1]Z\to Z$. Setting $\lambda_a=\xi_i=0$ represents
its zero section. Their composite on $Z$ is the identity.
\end{theorem}

\begin{proof}
The equations in \eqref{vc:cicp:eq:differential} are the partial
derivatives of $F$ with respect to the coordinates of $M$.
They therefore give its critical Koszul algebra.
They also present
$\operatorname{Sym}_{\mathcal O_Z}(\mathbb T_Z[1])$:
the shifted duals of the odd Koszul generators give the
degree-zero multipliers, and those of the base coordinates
give the $\xi_i$.

The common presentation identifies the standard shifted
cotangent form with the critical form. This is the local
comparison of Kinjo, Lemma~2.7.\footnote{T.~Kinjo,
\emph{Dimensional reduction in cohomological Donaldson--Thomas
theory}, \href{https://arxiv.org/abs/2102.01568v2}{arXiv:2102.01568v2},
Lemma~2.7. For the shifted cotangent form, see
D.~Calaque, \emph{Shifted cotangent stacks are shifted symplectic}.
\href{https://arxiv.org/abs/1612.08101v2}{arXiv:1612.08101v2},
Theorem~2.2.}

The two displayed algebra maps commute with the differential
and have the stated composite. The coordinate changes below
give their descent.
\end{proof}

\section{Descent and orientation}

The multipliers transform in the dual bundle. The odd base
coordinates must transform as well. Let $s'=G(x)s$ for an
invertible matrix of regular functions. Use column vectors for
$s,\lambda,\epsilon$ and put
\begin{equation}\label{vc:cicp:eq:frame}
 \lambda'=G^{-T}\lambda,\qquad \epsilon'=G\epsilon,
 \qquad
 \xi'_i=\xi_i+\lambda'^T(\partial_iG)\epsilon.
\end{equation}

\begin{proposition}\label{vc:cicp:prop:descent}
Equations~\eqref{vc:cicp:eq:frame}, together with ordinary cotangent
coordinate changes on $B$, glue the critical algebra and its
shifted symplectic form. Adding one coordinate $z$ and one
equation $s_{r+1}=z$ leaves the derived zero locus unchanged
and adds the hyperbolic term $\lambda_{r+1}z$ to its critical
potential.
\end{proposition}

\begin{proof}
The identity $\lambda^Ts=\lambda'^Ts'$ preserves $F$.
Moreover,
\[
 d\epsilon'=Gs=s',\qquad
 d\xi'_i=\lambda'^TG\partial_i s
             +\lambda'^T(\partial_iG)s
          =\lambda'^T\partial_i s'.
\]
Thus the correction in \eqref{vc:cicp:eq:frame} is exactly the
product-rule term required for a cochain map.

In the critical-chart convention the odd Liouville form is
\[
 \theta=\sum_i\xi_i\,d_{\mathrm{dR}}x_i
                +\sum_a\epsilon_a\,d_{\mathrm{dR}}\lambda_a.
\]
Since
$d_{\mathrm{dR}}\lambda'
 =G^{-T}d_{\mathrm{dR}}\lambda
  -G^{-T}(d_{\mathrm{dR}}G)^T\lambda'$, the extra term in
$\epsilon'^T d_{\mathrm{dR}}\lambda'$ cancels the correction
in $\sum_i\xi'_i d_{\mathrm{dR}}x_i$.
Hence $\theta$, and therefore its de Rham differential,
is preserved. Under two frame changes the correction satisfies
the product rule for their product matrix, proving the cocycle.

For the additional presentation pair, the four new generators
have differential
\[
 d\epsilon_{r+1}=z,\qquad d\xi_z=\lambda_{r+1}.
\]
Their algebra is the Koszul resolution of the origin in
$\mathbb A^2_{z,\lambda_{r+1}}$. It contracts to $\mathbb C$.
This argument does not justify deleting an original equation.
\end{proof}

An orientation of a critical locus is a square root of its
critical canonical line on the reduced support. A global
critical chart $\operatorname{Crit}(F)\subset M$ identifies
that line with $K_M^{\otimes2}$ on this support, so it
supplies the root $K_M$. If
$B\to H$ is smooth and the construction is relative to $H$,
the cotangent sequence of $M\to B$ gives
\begin{equation}\label{vc:cicp:eq:orientation}
 K_{M/H}\cong\pi^*(K_{B/H}\otimes\det V).
\end{equation}
Restricted to the reduced critical locus, its square is the
critical canonical line.

The corresponding absolute formula
uses $K_B$. These are orientations for the respective critical
charts. Identifying them with the intrinsic shifted-cotangent
orientation uses the determinant conventions of the comparison.

For the global spectral family, let $E$ be an elliptic curve,
$\Theta=\mathcal O_E(0)$ with its divisor section $\vartheta$,
and $B=E^{n-1}\times E_h$.
For an oriented graph edge $e$, let $u_e:B\to E$ be its
vertex difference after quotienting simultaneous translation.
Define
\[
 \mathcal S_e=u_e^*\Theta\otimes(u_e+h)^*\Theta,
 \qquad s_e=\vartheta(u_e)\vartheta(u_e+h),
 \qquad V=\bigoplus_e\mathcal S_e.
\]
The pairing with $s=(s_e)$ defines a global potential on
$\operatorname{Tot}_B(V^\vee)$.

Its relative critical locus uses derivatives along $B/E_h$
and the multipliers. The same local construction and descent
give the relative shifted cotangent completion.
The absolute critical locus also has an odd coordinate
$\xi_h$ with $d\xi_h=\partial_hF$ in a parameter chart.
These are two different enhancements of the family.

Formula~\eqref{vc:cicp:eq:orientation} uses
$\det V=\bigotimes_e\mathcal S_e$.
It needs no square root of an individual $\mathcal S_e$.
This critical orientation does not decide whether an individual
pair line has a root over the displacement base.
It does not decide whether the distinguished section is a square.

\section{Adding constraints}

Fix finitely many indexed line bundles $L_a$ over the same
smooth base $B$, with sections $s_a$. For a subset $I$ put
$V_I=\bigoplus_{a\in I}L_a$, and let $Z_I$ be the derived
zero locus of $(s_a)_{a\in I}$.
Let $\widehat A_I$ denote its critical algebra.
If $I\subset J$, forgetting equations defines
$f_{JI}:Z_J\to Z_I$. Cotangent pullback gives the correspondence
\begin{equation}\label{vc:cicp:eq:correspondence}
 T^*[-1]Z_I\ \longleftarrow\
 Z_J\mathbin{\times}_{Z_I}T^*[-1]Z_I
 \ \longrightarrow\ T^*[-1]Z_J,
\end{equation}
where all fiber products of derived schemes are derived.

\begin{theorem}\label{vc:cicp:thm:flags}
The correspondence \eqref{vc:cicp:eq:correspondence} is Lagrangian
for the difference of the two shifted symplectic forms.
In a common affine chart its middle algebra is
\begin{equation}\label{vc:cicp:eq:middle}
 W_{I,J}=\widehat A_I\otimes
          \Lambda(\epsilon_a:a\in J\setminus I),
 \qquad d\epsilon_a=s_a.
\end{equation}
For $I\subset J\subset K$, composition has the equivalence
\begin{equation}\label{vc:cicp:eq:composition}
 W_{I,J}\otimes^{\mathbb L}_{\widehat A_J}W_{J,K}
      \simeq W_{I,K}.
\end{equation}
These equivalences agree for every nested flag of subsets.
\end{theorem}

\begin{proof}
The left algebra map is inclusion from $\widehat A_I$.
The map from $\widehat A_J$ sets
$\lambda_a=0$ for $a\in J\setminus I$ and retains their
$\epsilon_a$. It fixes all other generators.
Equation~\eqref{vc:cicp:eq:differential} verifies that this is a
cochain map, including the differentials of the $\xi_i$.

The pulled-back critical forms agree because each extra
term contains $d_{\mathrm{dR}}\lambda_a=0$.
The correspondence is the shifted conormal of the graph of
$f_{JI}$. Its Lagrangian nondegeneracy is Calaque's shifted
conormal theorem, Theorem~2.8 of the cited paper.

As a module over $\widehat A_J$, the algebra $W_{J,K}$ has
a finite filtration by the exterior length of its new
generators, with free graded pieces. It is therefore
$K$-flat, so ordinary tensor product computes the derived
tensor product in \eqref{vc:cicp:eq:composition}.
It sets to zero precisely the multipliers in $K\setminus I$
and retains precisely their odd generators and equations.
This is $W_{I,K}$.

Every longer flag gives the same algebra.
Reordering the odd generators uses the exterior Koszul signs.
Frame changes respecting the indexed line summands preserve
these subset maps. Proposition~\ref{vc:cicp:prop:descent}
therefore glues the equivalences. An arbitrary change mixing
an included and an excluded equation changes the specified flag.
\end{proof}

This constructs geometric correspondences and their composition.
A map on their unintegrated vanishing-cycle coefficients still
requires a coefficient operation with the appropriate support
and orientation. The geometric composition alone does not
supply that operation.

\section{The completed triangle}

For a fixed complex parameter $h$, take the local shifted pair equations
\[
 s_1=x(x+h),\qquad s_2=y(y+h),\qquad
 s_3=(x+y)(x+y+h).
\]
They are the formal equations of the selected forward pair
sections $\vartheta(u)\vartheta(u+h)$.

For the two-arrow weights of Yang--Zhao, the original equal
equivariant parameter is $t$ and $h=2t$.\footnote{
Y.~Yang and G.~Zhao, \emph{Quiver varieties and elliptic
quantum groups}.
\href{https://arxiv.org/abs/1708.01418v1}{arXiv:1708.01418v1},
Remark~2.2(2) and equation~(3). Only these ordered intercolor
pair factors are used here.}

Integrating an invariant differential gives a holomorphic
local logarithm. The theta sections are their linear factors
times invertible holomorphic functions. Rescaling the dual
multipliers by those functions gives these equations and an
analytic isomorphism of their potential germs near zero.
The same changes may also be completed formally.

The invertible row change replacing $s_3$ by
$(s_3-s_1-s_2)/2=xy$ changes the full potential to
\begin{equation}\label{vc:cicp:eq:triangle-potential}
 W_h=P x(x+h)+Q y(y+h)+Txy,
 \qquad P=p+r,\quad Q=q+r,\quad T=2r.
\end{equation}
Thus its critical algebra has five odd generators and
\begin{align}
 d\epsilon_P&=x(x+h),& d\epsilon_Q&=y(y+h),&
 d\epsilon_T&=xy,\notag\\
 d\xi_x&=(2x+h)P+Ty,&
 d\xi_y&=(2y+h)Q+Tx.
 \label{vc:cicp:eq:triangle-gradient}
\end{align}

\begin{proposition}\label{vc:cicp:prop:triangle-strata}
For $h\ne0$, the reduced critical locus of $W_h$ is the
disjoint union of the three affine lines
\[
 (x,y;P,Q,T)=(0,0;0,0,T),\quad
 (-h,0;0,T,T),\quad (0,-h;T,0,T).
\]
The Hessian has rank four at every point of these lines.
For $h=0$, the reduced critical support is
$\mathbb A^3_{P,Q,T}$ at $x=y=0$.
Put $\Delta=T^2-4PQ$. Its tangent and obstruction dimensions are
\[
\begin{array}{c|c|c}
 \text{multiplier stratum}&\operatorname{rank}\operatorname{Hess}W_0
       &(\dim H^0\mathbb T,\dim H^1\mathbb T)\\\hline
 \Delta\ne0&2&(3,3)\\
 \Delta=0,\ (P,Q,T)\ne0&1&(4,4)\\
 (P,Q,T)=0&0&(5,5).
\end{array}
\]
\end{proposition}

\begin{proof}
The first three equations in \eqref{vc:cicp:eq:triangle-gradient}
give the three listed pairs $(x,y)$ when $h\ne0$.
The last two then give the listed multiplier lines.
The principal Hessian minor on $x,y,P,Q$ has determinant
$h^4$ on every line. Its tangent vector is in the Hessian
kernel, so the rank is exactly four.

At $h=0$, the equations $x^2=y^2=xy=0$ force $x=y=0$
on the reduced support. The only nonzero Hessian block is
\begin{equation}\label{vc:cicp:eq:normal-hessian}
 H_{\mathrm{normal}}=
 \begin{pmatrix}2P&T\\T&2Q\end{pmatrix},
 \qquad \det H_{\mathrm{normal}}=-\Delta.
\end{equation}
The critical tangent complex is the Hessian map between
two five-dimensional spaces. Its kernel and cokernel have
dimension five minus this rank.
\end{proof}

The completion has introduced noncompact multiplier directions.
Its critical tangent dimensions therefore differ from the
cohomology dimensions of the coherent incidence algebra.
Vanishing cycles on those multiplier directions provide a
third, explicitly computable coefficient object.
\section{Vanishing cycles and a single multiplier}
\label{vc:sec:vanishing}

For a complex algebraic variety $Y$, write
$\mathrm D_c^b(Y;\mathbb Q)$ for the bounded derived
category of rational sheaves on its analytic space, with
finite-dimensional algebraically constructible cohomology.
All constructible coefficients below lie in this category.

Write $\varphi_f=\phi_f[-1]$ for shifted vanishing
cycles, with $\varphi_0=\mathrm{id}$. On a smooth complex
$N$-fold $M$, put
\[
 \mathcal P_f=\varphi_f\mathbb Q_M[N].
\]
This is the perverse normalization. Its stalk at a critical
point $p$ is the reduced cohomology of a local Milnor fiber,
shifted so that
\begin{equation}\label{vc:vic:eq:milnor-shift}
 \mathcal H^k(\mathcal P_f)_p
       =\widetilde H^{k+N-1}(\operatorname{MF}_p(f);\mathbb Q).
\end{equation}
When Tate twists are displayed, coefficients are understood
in the corresponding mixed-Hodge or monodromic refinement.


The middle perverse convention places a local system on a
smooth complex $d$-fold in degree $-d$. The functor
$\varphi_f$ is exact for this perverse structure.
For an open inclusion $j$, the middle extension $j_{!*}K$
is the perverse image of
${}^pH^0(j_!K)\to{}^pH^0(Rj_*K)$.
Here ${}^pH^0$ is degree-zero perverse cohomology and $K$
is perverse on the open stratum. This extension has no
nonzero perverse subobject or quotient supported on the boundary.

Hyperbolic stabilization of a critical chart, with the
induced orientation, preserves this normalized coefficient.
The comparison includes the determinant double cover of the
added quadratic form.\footnote{C.~Brav, V.~Bussi, D.~Dupont,
D.~Joyce and B.~Szendr\H{o}i, \emph{Symmetries and stabilization
for sheaves of vanishing cycles}.
\href{https://arxiv.org/abs/1211.3259}{arXiv:1211.3259},
Theorem~5.4.} In the coordinate-equation stabilization above,
the specified hyperbolic pair gives that orientation.

We first justify the use of global fibers in the homogeneous
examples below. Let $g$ be a nonconstant homogeneous polynomial
of positive degree on $\mathbb C^N$ and suppose that
$g^{-1}(1)$ is smooth. Dilation identifies
\[
 g^{-1}(\delta)\cap B_\varepsilon
       \cong g^{-1}(1)\cap B_{\varepsilon|\delta|^{-1/\deg g}}.
\]

The squared norm on the smooth affine fiber is proper and
semialgebraic. Its critical-value set is semialgebraic and,
by Sard's theorem, has measure zero, hence is finite.
Above its last critical value, increasing the sublevel radius
does not change homotopy type. Thus the local Milnor fiber at
the origin has the homotopy type of $g^{-1}(1)$.
This argument allows a nonisolated critical locus.

Let $f_m(\lambda,z)=\lambda z^m$ on $\mathbb A^2$, for
$m\ge1$. For $m\ge2$ its reduced critical support is the
$\lambda$-axis. Let
$j:\mathbb C^*\hookrightarrow\mathbb C$ and
$i:\{0\}\hookrightarrow\mathbb C$ be its two strata.
On the punctured axis define the local system $L_m$ by the
reduced cyclic permutation representation
\begin{equation}\label{vc:vic:eq:cyclic-system}
 (L_m)_\lambda=
 \bigl\{(a_1,\ldots,a_m)\in\mathbb Q^m:
                         \textstyle\sum_j a_j=0\bigr\},
 \qquad T(a_1,\ldots,a_m)=(a_m,a_1,\ldots,a_{m-1}).
\end{equation}
Reversing the loop orientation replaces $T$ by its inverse.
For $m=1$ put $L_1=0$.

\begin{theorem}\label{vc:vic:thm:monomial}
On its reduced support, the underlying rational perverse sheaf
of $f_m$ is
\begin{equation}\label{vc:vic:eq:monomial-splitting}
 \mathcal P_{f_m}\cong j_{!*}L_m[1]\oplus i_*\mathbb Q.
\end{equation}
The extension is clean:
$j_!L_m[1]\simeq Rj_*L_m[1]$.
At the origin the only nonzero rational stalk group is
$\mathbb Q$ in degree zero.
\end{theorem}

\begin{proof}
For $\lambda\ne0$, the transverse fiber has $m$ points.
Their reduced degree-zero cohomology gives
\eqref{vc:vic:eq:cyclic-system}. The smooth axis supplies the
shift by one. The characteristic polynomial of $T$ on this
space is
\[
 1+t+\cdots+t^{m-1},\qquad \det(1-T)=m.
\]

Consequently the punctured-disc cochain complex
$[L_m\xrightarrow{T-1}L_m]$ is acyclic. This proves
$j_!L_m\simeq Rj_*L_m$. Both its stalk and costalk at the
origin vanish, so adjunction gives no morphisms, in any degree,
between this extension and an origin-supported complex.
The localization triangle therefore splits off the origin part.

The homogeneous fiber $\lambda z^m=1$ is $\mathbb C^*$,
parametrized by $z\ne0$. The preceding homogeneous argument
identifies the local Milnor fiber with it up to homotopy.
Its only reduced cohomology is
$H^1(\mathbb C^*;\mathbb Q)=\mathbb Q$.
Equation~\eqref{vc:vic:eq:milnor-shift} puts this in degree zero.
For $m=1$ the critical point is isolated and the same calculation
gives the entire sheaf. This proves
\eqref{vc:vic:eq:monomial-splitting} in rational coefficients.
\end{proof}

The circle complex also gives
$R\Gamma_c(\mathbb C^*,L_m)=0$, by duality or by the
same invertibility of $T-1$. Thus compact integration over
the multiplier retains the origin term alone.
The unintegrated local system still has rank $m-1$ and
cyclic monodromy. The coherent algebra $\mathbb C[z]/(z^m)$,
these vanishing cycles and their compact integral are three
different coefficient objects.

\section{The triangle stalks}

For the polynomial
$W_h=P x(x+h)+Q y(y+h)+Txy$ of
\eqref{vc:cicp:eq:triangle-potential}, write
$\mathcal P_h=\varphi_{W_h}\mathbb Q_{\mathbb A^5}[5]$.
At $h=0$ it is supported on $\mathbb A^3_{P,Q,T}$.
The discriminant $\Delta=T^2-4PQ$ defines the quadratic rank strata.

\begin{theorem}\label{vc:vic:thm:triangle-stalks}
At a point $(0,0;P,Q,T)$ the nonzero rational stalk groups of
$\mathcal P_0$ are as follows:
\begin{equation}\label{vc:vic:eq:stalk-table}
\begin{array}{c|c|c}
 \text{stratum}&\text{cohomological degree}&\text{rank}\\\hline
 \Delta\ne0&-3&1\\
 \Delta=0,\ (P,Q,T)\ne0&-2&1\\
 (P,Q,T)=0&-1&1.
\end{array}
\end{equation}
For $h\ne0$, its restriction to each of the three critical
lines is isomorphic to $\mathbb Q_{\mathbb A^1}[1]$.
\end{theorem}

\begin{proof}
At quadratic rank two, the holomorphic Morse lemma with
parameters gives a nondegenerate quadratic in two normal
variables and three smooth spectator coordinates.
Its Milnor fiber has reduced cohomology in degree one.
The shift in \eqref{vc:vic:eq:milnor-shift} therefore gives
$1-(5-1)=-3$.

At rank one, assume $P\ne0$. The patch with $Q\ne0$ is
identical after interchanging $x,y$. Completing the square gives
\[
 W_0=P\left(x+\frac{T}{2P}y\right)^2
              +\left(Q-\frac{T^2}{4P}\right)y^2.
\]
Choose a local square root of the unit $P$. The new coefficient
$s=Q-T^2/(4P)$ is a transverse parameter to the rank-one
stratum. The germ becomes $u^2+sv^2$, with two spectator
coordinates.

Theorem~\ref{vc:vic:thm:monomial} for $sv^2$
has reduced Milnor cohomology in degree one at its origin.
Thom--Sebastiani for the additional nondegenerate square
suspends it to degree two. Formula~\eqref{vc:vic:eq:milnor-shift}
then gives $-2$.

At the zero quadratic form, project the global fiber $W_0=1$
to $(x,y)$. The point $(0,0)$ is excluded. At every other
point, $P x^2+Txy+Qy^2=1$ is a nonzero affine linear
equation in three multipliers, so its solutions form an
affine plane. On $x\ne0$ one solves for $P$, and on $y\ne0$
one solves for $Q$. This is a locally trivial affine-bundle
torsor over $\mathbb C^2\setminus\{0\}$, with contractible
fiber. It has the homotopy type of $S^3$.

The homogeneous argument applies to $W_0$ of degree three,
so this is also the local Milnor homotopy type. Its unique
reduced class gives degree $3-(5-1)=-1$.

For $h\ne0$, Proposition~\ref{vc:cicp:prop:triangle-strata}
gives a smooth critical line and a nondegenerate Hessian
in four normal directions. The Morse lemma with parameters
gives a rank-one local system shifted by one on that line.
The line is simply connected, so this local system is constant.
\end{proof}

Stalk ranks alone do not describe how a perverse sheaf glues
across these strata. Even on the open quadratic stratum the
rank-one coefficient has nontrivial monodromy.

\begin{proposition}\label{vc:vic:prop:determinant-cover}
On $\Delta\ne0$, the normal vanishing cohomology of $W_0$ is the
sign local system associated with the double cover
\begin{equation}\label{vc:vic:eq:determinant-cover}
 w^2=\Delta.
\end{equation}
It is nontrivial.
\end{proposition}

\begin{proof}
A change of orthonormal frame for a nondegenerate complex
quadratic form acts on its one-dimensional vanishing
cohomology by the determinant character of $O(2,\mathbb C)$.
The two choices of determinant frame give the double cover
\eqref{vc:vic:eq:determinant-cover}. The determinant of the
normal Hessian is $-\Delta$ and multiplication of a square
root by $i$ identifies the two sign conventions.
For an explicit loop, set $Q=1,T=0$ and move $P$ once
around zero. A local normal coordinate $\sqrt P\,x$
returns with the opposite sign. Reflection of that coordinate
acts by $-1$ on the vanishing class.
\end{proof}

This double cover lies over multiplier space. It persists
at zero spectral displacement and is independent of the
cover $h=2t$ used to take roots of spectral pair lines.
It also coexists with the global critical orientation
\eqref{vc:cicp:eq:orientation}. An orientation and a constant
trivialization of every normal vanishing-cycle local system
are different data.

\section{The central perverse sheaf}
\label{vc:sec:central-sheaf}

Put $B_{\mathrm{mult}}=\mathbb A^3_{P,Q,T}$ and
$f=W_0=Px^2+Txy+Qy^2$. Regard
$\Phi=\mathcal P_0$ as a perverse sheaf on its reduced
support $B_{\mathrm{mult}}$.
Write $\mathcal T_f$ for raw geometric monodromy in the
value of $f$, with no dimension-dependent sign.
This is distinct from continuation in multiplier space
and from monodromy in the displacement $h$.
Let $D=\{\Delta=0\}$,
$D^\circ=D\setminus\{0\}$ and
$U=B_{\mathrm{mult}}\setminus D$.

\begin{lemma}\label{vc:cqm:lem:function-monodromy}
The automorphism $\mathcal T_f$ satisfies $\mathcal T_f^2=1$.
The only nonzero stalk cohomology on each stratum is
\begin{equation}\label{vc:cqm:eq:stalks}
\begin{array}{c|c|c}
 \text{stratum}&\text{stalk of }\Phi&\mathcal T_f\\ \hline
 U&H^{-3}=\mathbb Q&+1\\
 D^\circ&H^{-2}=\mathbb Q&-1\\
 \{0\}&H^{-1}=\mathbb Q&+1.
\end{array}
\end{equation}
\end{lemma}

\begin{proof}
The path $(x,y)\mapsto(e^{\pi it}x,e^{\pi it}y)$, with
$P,Q,T$ fixed, rotates the value of $f$ once. It preserves
small balls centered on $B_{\mathrm{mult}}$. Its endpoint is the involution
$(x,y)\mapsto(-x,-y)$, and two circuits give the identity.
Thus it defines the stated involution on vanishing cycles.

The local fibers are those of
Theorem~\ref{vc:vic:thm:triangle-stalks}. At rank two, a
nondegenerate binary quadratic form has Milnor fiber of
homotopy type $S^1$. At rank one, completing a square gives
$u^2+sv^2$, with two additional smooth parameters.
The fiber of $sv^2$ has homotopy type $\mathbb C^*$.
Adding $u^2$ suspends its reduced cohomology. Its function
monodromy on that cohomology is the product of $+1$ for
$sv^2$ and $-1$ for $u^2$.

These are the Thom--Sebastiani monodromy conventions.\footnote{
C.~Brav, V.~Bussi, D.~Dupont, D.~Joyce and B.~Szendr\H{o}i,
\emph{Symmetries and stabilization for sheaves of vanishing
cycles}. \href{https://arxiv.org/abs/1211.3259}{arXiv:1211.3259},
Definition~2.12 and Theorem~2.13. Their twisted monodromy
includes the sign $(-1)^{\dim U}$. The symbol $\mathcal T_f$ here is
their untwisted function monodromy.}

At the origin, projection of $f^{-1}(1)$ to
$\mathbb C^2\setminus\{0\}$ is an affine-plane torsor:
for each nonzero $(x,y)$, the equation is a nonzero linear
equation in $P,Q,T$. It is locally trivial and has
contractible fibers. Hence $f^{-1}(1)\simeq S^3$.
The homogeneous comparison proved in
Section~\ref{vc:sec:vanishing} applies to this cubic, whose
fiber over one is smooth.

The involution covers the antipodal map of
$\mathbb C^2\setminus\{0\}$ and acts by $+1$ on $H^3$.
At rank two its action on $H^1(S^1)$ is also $+1$.
The shift $\phi_f[-1]\mathbb Q[5]$ places reduced Milnor
cohomology of degree $k$ in stalk degree $k-4$.
\end{proof}

\subsection{The two local systems}

Consider the double covers
\begin{equation}\label{vc:cqm:eq:covers}
 \nu:\widetilde B_{\mathrm{mult}}=\{w^2=\Delta\}\longrightarrow B_{\mathrm{mult}},
 \qquad
 \rho:\mathbb A^2\longrightarrow D,
 \quad(a,b)\longmapsto(a^2,b^2,2ab).
\end{equation}
Let $L_{\det}$ be the sign local system of $\nu$ on $U$,
and let $L_{\mathrm{sq}}$ be the sign local system of
$\rho$ on $D^\circ$. 
The rank-two vanishing line is $L_{\det}$ by
Proposition~\ref{vc:vic:prop:determinant-cover}.
On $D^\circ$, the form
is $(ax+by)^2$. The loop $P=e^{2\pi it}$, $Q=T=0$ reverses
its chosen square coordinate and hence its suspended
vanishing class. Since
$D^\circ=(\mathbb C^2\setminus\{0\})/\{\pm1\}$ has
fundamental group $\mathbb Z/2$, this identifies the
rank-one vanishing line with $L_{\mathrm{sq}}$.

Write $j:U\hookrightarrow B_{\mathrm{mult}}$, $k:D^\circ\hookrightarrow D$
and $i:D\hookrightarrow B_{\mathrm{mult}}$. For a pure-dimensional variety $Y$, the intersection
complex $\operatorname{IC}_Y$ is the middle extension of
the constant complex on its smooth locus, normalized by
$\operatorname{IC}_Y|_{Y_{\mathrm{sm}}}
=\mathbb Q[\dim Y]$.
For a finite double-cover pushforward, the superscript
$-$ denotes its deck anti-invariant direct summand.

\subsection{The vertex and its involution}

\begin{lemma}\label{vc:cqm:lem:node}
The stalk of $\operatorname{IC}_{\widetilde B_{\mathrm{mult}}}$ at the
origin has a copy of $\mathbb Q$ in each of degrees
$-3$ and $-1$. Deck monodromy acts by $+1$ and $-1$
on them, respectively.
\end{lemma}

\begin{proof}
The matrix
\[
 \begin{pmatrix}2P&T+w\\T-w&2Q\end{pmatrix}
\]
has determinant $w^2-\Delta$. It identifies $\widetilde B_{\mathrm{mult}}$
with the rank-at-most-one locus of two-by-two matrices.
Choosing the line containing the image gives a proper
resolution
\[
 r:\operatorname{Tot}_{\mathbb P^1}
       (\mathcal O(-1)\oplus\mathcal O(-1))
       \longrightarrow\widetilde B_{\mathrm{mult}}.
\]
It is an isomorphism away from the origin, with exceptional
fiber $\mathbb P^1$.

Since $2\dim\mathbb P^1=2<3$, the
map is small, and
$Rr_*\mathbb Q[3]=\operatorname{IC}_{\widetilde B_{\mathrm{mult}}}$.
Proper base change gives the two stalk groups.\footnote{
M.~A.~de Cataldo and L.~Migliorini,
\emph{The decomposition theorem, perverse sheaves and the
topology of algebraic maps}.
\href{https://arxiv.org/abs/0712.0349}{arXiv:0712.0349},
Remark~4.2.4.}

The deck involution is matrix transpose. It exchanges
the two small resolutions, so its action must be
computed intrinsically. The link of the node is the
circle bundle of $\mathcal O(-1,-1)$ over
$\mathbb P^1\times\mathbb P^1$. If $a,b$ are the two
degree-two generators, its Euler class is $a+b$ up to
sign. The Gysin sequence gives
\[
 H^2(\mathrm{link};\mathbb Q)
    =\frac{\mathbb Q a\oplus\mathbb Q b}
            {\mathbb Q(a+b)}.
\]

Transpose interchanges the two factors and acts by $-1$
on this quotient. Intersection extension across an
isolated threefold singularity retains this link group
as stalk degree $-1$, and the constant link group as
degree $-3$. The latter is invariant.
\end{proof}

\begin{theorem}\label{vc:cqm:thm:coefficient}
There is a decomposition of rational perverse sheaves
\begin{equation}\label{vc:cqm:eq:decomposition}
 \Phi=\mathcal P_+\oplus\mathcal P_-,
 \qquad \mathcal P_\pm=\operatorname{im}
                  \frac{1\pm\mathcal T_f}{2},
\end{equation}
with
\begin{align}
 \mathcal P_+&\cong j_{!*}(L_{\det}[3])
       \cong(R\nu_*\operatorname{IC}_{\widetilde B_{\mathrm{mult}}})^- ,
       \label{vc:cqm:eq:plus}\\
 \mathcal P_-&\cong i_*k_{!*}(L_{\mathrm{sq}}[2])
       \cong i_*(R\rho_*\mathbb Q_{\mathbb A^2}[2])^- .
       \label{vc:cqm:eq:minus}
\end{align}
Raw function monodromy acts by $+1$ and $-1$ on these
summands, respectively.
\end{theorem}

\begin{proof}
Lemma~\ref{vc:cqm:lem:function-monodromy} gives the two
projectors in the perverse category. Vanishing cycles
commute with Verdier duality, which replaces monodromy
by its inverse. Each summand is therefore self-dual.

The plus summand restricts to $L_{\det}[3]$ on $U$ and
has zero stalks and costalks on $D^\circ$. At the origin,
its stalk has degree $-1$ and its costalk degree $+1$.
It has no perverse quotient or subobject supported on the
boundary. These are precisely the middle-extension
conditions.

The minus summand has support $D$, restricts
to $L_{\mathrm{sq}}[2]$, and has zero stalk and costalk
at the origin. The same criterion identifies it.\footnote{
L.~G.~Maxim and J.~Sch\"urmann, \emph{Constructible sheaf
complexes in complex geometry and applications}.
\href{https://arxiv.org/abs/2105.13069}{arXiv:2105.13069},
Sections~2.2, 2.4 and 4.2, for intermediate extension,
finite pushforward and Verdier duality.}

Finite pushforward is perverse exact. For $\rho$, the
origin is a fixed point of the deck involution, which
acts trivially on the constant stalk. Its anti-invariant
stalk and, by duality, costalk vanish. This proves
\eqref{vc:cqm:eq:minus}.

For $\nu$, the anti-invariant stalk
vanishes at every smooth branch point. Its origin stalk
is $\mathbb Q$ in degree $-1$, with dual costalk in degree
$+1$, by Lemma~\ref{vc:cqm:lem:node}. Thus its
anti-invariant pushforward satisfies the same
middle-extension conditions as $\mathcal P_+$.
\end{proof}

\section{Dimensional reduction and constraint maps}
\label{vc:sec:reduction}

Return to a smooth complex $m$-fold $B$, a rank-$r$ bundle
$V$, and a section $s$. Put
$M=\operatorname{Tot}_B(V^\vee)$, let $\pi:M\to B$ be
its projection, and let $F=\langle\lambda,s\rangle$.
Write $i:Z_{\mathrm{cl}}\hookrightarrow B$ for the ordinary
zero locus with its underlying complex topology, and
$\widetilde i:M|_{Z_{\mathrm{cl}}}\hookrightarrow M$.
The symbol $R\pi_!$ means derived direct image with proper
support. No properness of the multiplier bundle is assumed.

There is a natural coefficient map
\begin{equation}\label{vc:eq:reduction-map}
\begin{aligned}
 \alpha_{\mathcal G}:R\pi_!\varphi_F\pi^*\mathcal G
 &\longrightarrow
 R\pi_!\widetilde i_*\widetilde i^*\pi^*\mathcal G\\
 &\cong i_*i^*\mathcal G[-2r](-r),
 \qquad \mathcal G\in\mathrm D_c^b(B;\mathbb Q).
\end{aligned}
\end{equation}
To specify the first arrow, let $i_0:F^{-1}(0)\hookrightarrow M$.
The support description
\[
 \varphi_F\mathcal H
 =i_{0*}i_0^*R\Gamma_{\{\operatorname{Re}F\leq0\}}\mathcal H
\]
gives its natural map to coefficients on the zero fiber.
Restrict this map to $M|_{Z_{\mathrm{cl}}}$, then apply
$R\pi_!$. The second arrow in \eqref{vc:eq:reduction-map}
uses the complex orientation of the rank-$r$ fibers.

\begin{theorem}\label{vc:prop:linear-reduction}
The map $\alpha_{\mathcal G}$ is an isomorphism, natural in
$\mathcal G$. Thus
\begin{equation}\label{vc:vic:eq:dimensional-reduction}
 R\pi_!\varphi_F\pi^*\mathcal G
 \simeq i_*i^*\mathcal G[-2r](-r).
\end{equation}
For rational sheaves the Tate twist is omitted.
The comparison commutes with raw function monodromy,
which is the identity on the right-hand side.
For two linear potentials it carries the Thom--Sebastiani
map to the compactly supported K\"unneth map.
\end{theorem}

\begin{proof}
Where a component of $s$ is nonzero, replace its multiplier
coordinate by $F$. A coefficient pulled back from $B$ is
constant in this coordinate, so its vanishing cycles are
zero. Thus the source is supported over $Z_{\mathrm{cl}}$.
On a trivializing chart, the linear-potential reduction
theorem identifies the map just defined with an
isomorphism.\footnote{B.~Davison, \emph{The critical CoHA
of a quiver with potential}.
\href{https://arxiv.org/abs/1311.7172}{arXiv:1311.7172},
Theorem~A.1, Corollary~A.7 and Proposition~A.8.}

The local cancellation can be seen in one multiplier.
Homogeneity makes the nearby term a derived direct image
from its punctured multiplier line. For
$j:\mathbb C^*\hookrightarrow\mathbb C$ and a local
system $L$ with fiber $A$ and monodromy $T$, both
$R\Gamma(\mathbb C^*,L)$ and $(Rj_*L)_0$ are the circle
complex $[A\xrightarrow{T-1}A]$.
Their restriction map is an isomorphism.
The localization triangle gives
$R\Gamma(\mathbb C,j_!L)=0$, and Verdier duality gives
$R\Gamma_c(\mathbb C,Rj_*L)=0$.
This also holds for bounded complexes with locally constant
cohomology and requires no finite-order hypothesis on $T$.

For several multipliers, the blow-up of the zero section
is the tautological line over the projective multiplier
bundle. On its $a$th chart the potential is
$\lambda(s_a+\sum_{b\ne a}s_bv_b)$.
The same punctured-line cancellation, with proper
base change for the blow-up, is the local reduction
argument. The remaining coefficient is
$i_*i^*\mathcal G\otimes R\Gamma_c(\mathbb A^r,\mathbb Q)$,
and the fiber complex is $\mathbb Q[-2r](-r)$.

The support map is invariant under bundle frames, so
these local comparisons glue. Its monodromic refinement
has constant value-parameter monodromy on the reduced
side. The support inclusion defining Thom--Sebastiani
commutes with restriction, giving the last assertion.
For $r=0$, the formula is $\varphi_0=\mathrm{id}$.
\end{proof}

The pullback condition is essential. For $F=t$ on
$\mathbb A^1$ and the origin-supported coefficient
$\mathcal H=i_{0*}\mathbb Q$, nearby cycles vanish and
$\varphi_t\mathcal H=\mathcal H$.
Its compact integral is $\mathbb Q$, whereas the
coefficient section $1$ over a point has empty zero locus.
Replacing $\pi^*\mathcal G$ by this arbitrary coefficient
would therefore make \eqref{vc:vic:eq:dimensional-reduction} false.

\begin{corollary}\label{vc:vic:cor:integrated-constraints}
For the subset $I$ of a fixed collection of sections, let
$r_I=|I|$, let $\pi_I$ be its multiplier bundle and let
$i_I:(Z_I)_{\mathrm{cl}}\hookrightarrow B$.
The normalized integrated coefficient
\begin{equation}\label{vc:vic:eq:normalized-coefficient}
 \mathcal C_I=
 R\pi_{I!}\varphi_{F_I}\mathbb Q[2r_I](r_I)
       \simeq i_{I*}\mathbb Q_{(Z_I)_{\mathrm{cl}}}
\end{equation}
has restriction maps $\mathcal C_I\to\mathcal C_J$ for
$I\subset J$, compatible with every nested flag.
\end{corollary}

\begin{proof}
The inclusion
$a:(Z_J)_{\mathrm{cl}}\hookrightarrow(Z_I)_{\mathrm{cl}}$
has the ordinary coefficient unit
\[
 \mathbb Q_{(Z_I)_{\mathrm{cl}}}
       \longrightarrow Ra_*\mathbb Q_{(Z_J)_{\mathrm{cl}}}.
\]
Push it to $B$ and transport it through
\eqref{vc:vic:eq:normalized-coefficient}.
The units for successive inclusions compose to the unit
for their composite. This proves the flag identities.
The normalization removes the separate shift and twist
for every added multiplier.
\end{proof}

These are maps after multiplier integration. They are not
unintegrated vanishing-cycle Gysin maps on a sheaf-extension
stack. In particular, they do not equip the geometric
correspondences of Theorem~\ref{vc:cicp:thm:flags} with a Hall
product before integration.

The right side of \eqref{vc:vic:eq:dimensional-reduction} sees
the underlying classical space. It retains neither nilpotent
thickening nor the negative coherent Koszul cohomology as such.
For $s=z^m$, this agrees with
Theorem~\ref{vc:vic:thm:monomial}: the integral is independent
of $m$, although the transverse local system and coherent
length both depend on $m$.

For $f_m=\lambda z^m$, the perverse normalization is
$[2]$. Dimensional reduction therefore makes its compact
integral $\mathbb Q(-1)$ in degree zero in the monodromic
mixed-Hodge refinement. The clean punctured local system
has zero compact integral. Consequently the origin term
of Theorem~\ref{vc:vic:thm:monomial} has that same Tate class.
This specifies the Hodge class without identifying local
and global mixed-Hodge structures merely from a homotopy
equivalence of their fibers.

For a fixed triangle parameter, $m=2$, $r=3$, and the
perverse normalization is $[5]$. Hence
\begin{equation}\label{vc:vic:eq:triangle-integral}
 R\Gamma_c(\operatorname{Crit}(W_h),\mathcal P_h)
       \simeq R\Gamma_c((Z_h)_{\mathrm{cl}},\mathbb Q)[-1](-3).
\end{equation}
At $h=0$ the right side has one class, and at a generic
nonzero displacement it has three classes, all in degree one.

For the global elliptic pair sections, take
$B=E_x\times E_y$ and the three differences $x,y,x+y$.
The first two pair sections give $x,y\in\{0,-h\}$ when
$h\ne0$. Three of the four pairs satisfy the last equation
for every $h$. At $(-h,-h)$ that section has value
$\vartheta(-2h)\vartheta(-h)$. It vanishes precisely at
nonzero two-torsion displacements. Thus the global version
of \eqref{vc:vic:eq:triangle-integral} has four classes there.

The compactness of $E^2$ does not make its multiplier
bundle proper. The symbol $!$ in these formulas remains
essential. Verdier duality converts it to an ordinary direct
image only while changing the coefficients to their duals.

\section{Compact integration and monodromy}

\begin{proposition}\label{vc:cqm:prop:integral}
Compact integration over multiplier space satisfies
\begin{equation}\label{vc:cqm:eq:integrals}
 R\Gamma_c(B_{\mathrm{mult}},\mathcal P_-)=0,\qquad
 R\Gamma_c(B_{\mathrm{mult}},\mathcal P_+)\cong\mathbb Q[-1].
\end{equation}
In particular, it annihilates a nonzero perverse sheaf.
\end{proposition}

\begin{proof}
By \eqref{vc:cqm:eq:minus}, the first integral is the
anti-invariant part of
$R\Gamma_c(\mathbb C^2,\mathbb Q)[2]$.
The only compact cohomology of $\mathbb C^2$ is
$H_c^4=\mathbb Q$, and $(a,b)\mapsto(-a,-b)$ preserves
its complex orientation. That anti-invariant part is zero.

For the plus summand, the resolution in
Lemma~\ref{vc:cqm:lem:node} retracts to its exceptional
$\mathbb P^1$. Ordinary global cohomology of its
pushforward therefore maps isomorphically to the
origin stalk. This restriction is intrinsic and
commutes with the deck involution on the intersection
complex. No action on a chosen resolution is required.
Consequently
$R\Gamma(B_{\mathrm{mult}},\mathcal P_+)=\mathbb Q[1]$.
Verdier duality and self-duality give the second formula.
\end{proof}

These statements are in rational constructible sheaves.
For the total complex $\Phi$, dimensional reduction gives
$\mathbb Q[-1](-3)$ with its monodromic mixed Hodge
normalization, as in \eqref{vc:vic:eq:triangle-integral} and
Theorem~\ref{vc:prop:linear-reduction}.
This does not specify separate Hodge lifts of
\eqref{vc:cqm:eq:plus} and \eqref{vc:cqm:eq:minus}.

The same reduction has a useful monodromy consequence.
Let $X$ be smooth, $s_1,\ldots,s_r$ regular functions,
$\pi:X\times\mathbb A^r\to X$, and
$F=\sum_{a=1}^r\lambda_a s_a$.

\begin{proposition}\label{vc:cqm:prop:linear-monodromy}
Let $\mathcal K$ be a monodromy-invariant direct
summand of $\phi_F[-1]\mathbb Q$ on which
$\mathcal T_F-1$ is invertible. Then $\pi_!\mathcal K=0$.
\end{proposition}

\begin{proof}
The natural reduction map identifies
$\pi_!\phi_F[-1]\mathbb Q$ with the constant
coefficient on $Z=\{s_a=0\}$, shifted by $-2r$.
It commutes with raw function monodromy, and the
target monodromy is the identity. Hence on the
direct summand $\pi_!\mathcal K$ the endomorphism
$\mathcal T_F-1$ is both zero and invertible.

The required equivariance is part of the natural
monodromic reduction map.\footnote{B.~Davison,
\emph{The critical CoHA of a quiver with potential}.
\href{https://arxiv.org/abs/1311.7172}{arXiv:1311.7172},
Theorem~A.1 and Corollary~A.7. In the proof of the latter,
the comparison over the punctured value parameter is
constant on the reduced side.}
\end{proof}

The pullback hypothesis and its counterexample are those
of Theorem~\ref{vc:prop:linear-reduction}.

\subsection{The integral cyclic boundary}

Over $\mathbb Z$, division by two is unavailable.
The cyclic local system of the monomial example gives
a direct test of the resulting boundary terms.
For $m\ge2$, let
\[
 I_m=\{(a_1,\ldots,a_m)\in\mathbb Z^m:
                         \textstyle\sum a_i=0\},
\]
with cyclic permutation $T_m$.

\begin{proposition}\label{vc:cqm:prop:integral-link}
The map $T_m-1:I_m\to I_m$ is injective and has
cokernel $\mathbb Z/m$.
\end{proposition}

\begin{proof}
Let $C_m$ be the cyclic group of order $m$.
Multiplication by $t-1$ identifies
$\mathbb Z[C_m]/\mathbb Z(1+\cdots+t^{m-1})$ with
its augmentation ideal $I_m$.
The invariant vectors in $I_m$ are constant vectors
with sum zero, hence vanish. Taking coinvariants in
the quotient presentation gives
$\mathbb Z/(1+\cdots+1)=\mathbb Z/m$.
Equivalently, the Smith factors are
$1,\ldots,1,m$.
\end{proof}

For the associated local system on a punctured disk,
the boundary stalk of its derived direct image is the
circle complex $[I_m\xrightarrow{T_m-1}I_m]$ in
degrees $0,1$. It is nonzero over $\mathbb Z$ and
vanishes after tensoring with $\mathbb Q$.
In particular, a sign local system has an integral
$\mathbb Z/2$ boundary term. This computation does not
give an integral direct-sum version of
Theorem~\ref{vc:cqm:thm:coefficient}.
\section{Three branches specializing to one point}

Work over a small complex disk $\mathbb D$ around $h=0$.
The ordinary total zero locus of the polynomial triangle
has ring
\begin{equation}\label{vc:vic:eq:three-branch-ring}
 \mathbb C[h,x,y]/(x(x+h),y(y+h),xy).
\end{equation}
Define
\[
 z_0=h+x+y,\qquad z_1=-x,\qquad z_2=-y.
\]
Then $h=z_0+z_1+z_2$, and its ideal is
$(z_0z_1,z_0z_2,z_1z_2)$.
Thus the total classical space is reduced: it is the union
of three coordinate axes. Each axis maps isomorphically to
the parameter disc and all three meet at its origin.

The nonreduced central scheme fiber has just one underlying
topological point.

Let $f:Z_{\mathrm{cl}}\to\mathbb D$ be this finite map and put
$\mathcal V=f_*\mathbb Q_{Z_{\mathrm{cl}}}$.
Write $i:\{0\}\hookrightarrow\mathbb D$ for the origin inclusion.

\begin{theorem}\label{vc:vic:thm:specialization}
There is an exact sequence of constructible sheaves
\begin{equation}\label{vc:vic:eq:normalization-sequence}
 0\longrightarrow\mathcal V\longrightarrow\mathbb Q_{\mathbb D}^3
   \xrightarrow{\ (a,b,c)\mapsto(b-a,c-a)\ }
      i_*\mathbb Q^2\longrightarrow0.
\end{equation}
The nearby cycles at zero are $\mathbb Q^3$ with identity
monodromy. The specialization map is
\begin{equation}\label{vc:vic:eq:specialization}
 \mathbb Q\longrightarrow\mathbb Q^3,
       \qquad a\longmapsto(a,a,a),
\end{equation}
and consequently
\begin{equation}\label{vc:vic:eq:parameter-cycles}
 \varphi_h(\mathcal V[1])_0\cong\mathbb Q^2.
\end{equation}
Branch permutations induce the standard two-dimensional
representation of $S_3$ on this last vector space.
There is no nonzero subsheaf of $\mathcal V$ supported only
at the origin.
\end{theorem}

\begin{proof}
Normalization separates the three axes into three discs.
A locally constant function descends exactly when its
three values at the common point agree. On the punctured
disc normalization is an isomorphism on each branch. At
the origin its coefficient map is the diagonal inclusion.
These stalk calculations prove
\eqref{vc:vic:eq:normalization-sequence} and
\eqref{vc:vic:eq:specialization}.

Each branch is a section over the whole disc, so continuation
around zero fixes it. The vanishing-cycle triangle therefore
identifies \eqref{vc:vic:eq:parameter-cycles} with
$\mathbb Q^3/\mathbb Q(1,1,1)$.
Permutation of the three branch values gives the stated
$S_3$ representation.
Finally, a section supported only at the common point must
vanish along every incident branch. Its common value is
then zero. This proves the last assertion.
\end{proof}

In particular, $\mathcal V$ cannot be replaced by a punctured
constant sheaf together with an independent central
skyscraper while preserving its specialization map.
Identity nearby monodromy has not made specialization an
isomorphism.

Let $p:(x,y,h)\mapsto h$ be the base projection, so that
$f=p|_{Z_{\mathrm{cl}}}$. Let $\pi$ project the multiplier
bundle to this three-dimensional base.
For the absolute total-family potential $F$ on that
six-dimensional bundle, dimensional reduction before
parameter restriction gives
\begin{equation}\label{vc:vic:eq:absolute-family}
 R(p\pi)_!\varphi_F\mathbb Q\simeq\mathcal V[-6](-3).
\end{equation}

Indeed $f$ is finite, and proper base change to its finite
fibers shows that its higher constructible direct images
vanish. Equation~\eqref{vc:vic:eq:dimensional-reduction} then
gives \eqref{vc:vic:eq:absolute-family}.
This uses the absolute potential with its $h$ derivative.
It asserts no arbitrary base-change identity for unintegrated
vanishing cycles. The fixed-$h$ perverse coefficients have
the separate normalization in \eqref{vc:vic:eq:triangle-integral}.

The coherent family and the constructible family therefore
have distinct defects. Monic reduction of the first two
Koszul equations gives
\[
 A_h=\mathbb C[[h]][x,y]/(x(x+h),y(y+h))
            \otimes\Lambda(\eta),\qquad d\eta=xy.
\]
Multiplication by $xy$ on the basis $1,x,y,xy$ has sole
nonzero row $(1,-h,-h,h^2)$. Its unit entry gives free
cohomology of rank three in each of degrees $0,-1$.

At zero it is the nonformal algebra of
Theorem~\ref{ei:gio:thm:triangle-minimal}.
Meanwhile the constructible coefficient has the rank-two
specialization defect \eqref{vc:vic:eq:parameter-cycles}.
These statements coexist because their coefficient functors
retain different data.

There is a concrete evaluation map between the coherent
degree-zero ring and its branch normalization. For
$R=\mathbb C[[h]]$, evaluate the basis $1,x,y$ at
$(0,0),(-h,0),(0,-h)$. The resulting matrix is
\begin{equation}\label{vc:vic:eq:coherent-evaluation}
 E_h=\begin{pmatrix}1&0&0\\1&-h&0\\1&0&-h\end{pmatrix},
 \qquad \det E_h=h^2.
\end{equation}
It is injective over $R$, and its image consists exactly of
triples whose pairwise differences are divisible by $h$.
Indeed subtracting the first row from the other two makes
the matrix $\operatorname{diag}(1,-h,-h)$.

Thus its cokernel is $(R/(h))^2$, and its specialization
has rank one. After complexifying the rational stalks,
the surviving diagonal is the one in
\eqref{vc:vic:eq:specialization}. This compares the displayed
evaluation maps. It does not identify coherent and
constructible coefficient functors or their higher products.

Near a nonzero two-torsion displacement, the fourth elliptic
tree branch meets the last divisor in a simple isolated point,
disjoint from the three persistent branches.
Its constructible contribution is an actual parameter-supported
skyscraper. The corresponding coherent local complex is
$[\mathcal O\xrightarrow{v}\mathcal O]$, where $v$ is a
parameter times a unit. Its extra negative derived-fiber
class comes from Tor. This isolated contribution differs
from the collision of the three persistent branches at zero.

\section{The compact coefficient comparison}
\label{vc:sec:compact}

The surface Hall theory supplies an independent comparison
source. For a projective K3 surface $S$, let
$\mathfrak M_S(\gamma)$ be the stack of coherent sheaves
of Mukai vector $\gamma$. Write $v_\gamma$ for its
virtual dimension. Borel--Moore homology is defined by
the dualizing complex, with proper direct image as its
covariant operation.

The shifted rational groups
\[
 H_{S,n}(\gamma)
 =H^{\mathrm{BM}}_{-n+v_\gamma}
       (\mathfrak M_S(\gamma);\mathbb Q)
\]
carry the Hall multiplication and the cocommutative
coproduct constructed by Davison--Hennecart--Kinjo--Schiffmann--Vasserot.

The coproduct takes values in the product
completion over charge decompositions, while multiplication
uses the locally finite part. Let $\chi(\gamma_1,\gamma_2)$
be the Euler pairing, the alternating dimension of the
surface Ext groups. The charge braiding includes
$(-1)^{\chi(\gamma_1,\gamma_2)}$, together with the
cochain Koszul sign. The primitives for this coproduct
are their affinized BPS Lie algebra.\footnote{
B.~Davison, L.~Hennecart, T.~Kinjo, O.~Schiffmann and
E.~Vasserot, \emph{Hecke operators on symplectic surfaces
and $\chi$-independence}.
\href{https://arxiv.org/abs/2607.01355v1}{arXiv:2607.01355v1},
equation~(4.1), Section~6.1, Proposition~6.4,
Corollary~6.5 and Definition~6.6.}

Relative dimensional reduction also exists for perverse
pullbacks along locally finitely presented geometric
morphisms of derived Artin stacks, with classical base
locally of finite type.\footnote{A.~A.~Khan, T.~Kinjo,
H.~Park and P.~Safronov, \emph{Period sheaves via perverse
pullbacks}.
\href{https://arxiv.org/abs/2510.18610v1}{arXiv:2510.18610v1},
Theorem~4.1. The dimension is the relative virtual dimension
of the derived source.}
For a derived zero locus of rank $r$ on a smooth
$m$-dimensional base, this dimension is $m-r$.
It agrees with the shift in
\eqref{vc:vic:eq:triangle-integral} for $m=2,r=3$.

An application to a compact translation quotient must
identify its coefficient and extension maps with those
in the theorem. An Ext-quiver comparison must also
intertwine the surface coproduct before any primitive
identification follows.

\begin{question}\label{vc:vic:qn:compact-comparison}
For the common-torsor extension geometry of sheaves on
$S\times E_{\mathrm{geom}}$, construct the reduced oriented
critical coefficient functor and its maps on extension
correspondences. Give the map to a spectral base.
Specify its recipient: a coherent incidence algebra,
an unintegrated vanishing-cycle coefficient, its
constructible reduction, or an independently defined
state complex.

Can this comparison preserve the common translation torsor,
relative empty object, orientations, supports and every
flag composition? Determine which higher products,
monodromy summands and specialization maps survive it.
Construct its marked state and compare compact reduced
integration with the selected spectral trace.
Finally, identify the geometric modular line and section
required by the Igusa denominator.
\end{question}
